\section{Marc teòric}

\subsection{Evolució històrica d'Honda}
\textit{[Desenvolupament de l'origen de l'empresa, evolució tecnològica i posicionament mundial...]}

\subsection{Classificació de motocicletes}

\subsubsection{Motocicletes nues}
Les motocicletes nues (*naked*) es caracteritzen per l'absència total o parcial de carenat, una posició de conducció vertical i un estil minimalista. Estan dissenyades per a un ús diari amb un equilibri entre rendiment i potència.
\textit{Exemples: Honda CB650R, CB500 Hornet, CB750 Hornet.}

\subsubsection{Motocicletes esportives}
Les motocicletes esportives (*super sport*) disposen de carenat per reduir la resistència de l'aire i augmentar l'aerodinàmica. Estan optimitzades per al circuit i la velocitat, amb una posició de conducció agressiva i inclinada.
\textit{Exemples: Honda CBR600RR, Kawasaki Ninja ZX-6R, Yamaha YZF-R6.}

\subsubsection{Motocicletes d'aventura}
Dissenyades per a un ús mixt (asfalt i terra), les motocicletes d'aventura (*adventure* o *trail*) utilitzen pneumàtics mixtos i suspensions de llarg recorregut (entre 150 i 200 mm).
\textit{Exemples: Honda Africa Twin, Honda XL750 Transalp, Yamaha Ténéré 700.}

\subsubsection{Motocicletes de turisme}
Les motocicletes de turisme (*touring*) estan enfocades a viatges de llarga distància. Disposen de carenats extensos, gran capacitat de càrrega amb maletes i seients molt ergonòmics.
\textit{Exemples: Honda NT1100, Honda GL1800 Gold Wing Tour, BMW R1250 RT.}

\subsubsection{Motocicletes de turisme esportiu}
Les motocicletes de turisme esportiu (*sport touring*) són un híbrid que manté la capacitat de viatge del segment anterior però amb una part de cicle i un motor més potents i àgils.
\textit{Exemples: Suzuki GSX-S 1000 GT, Kawasaki Ninja 1000 SX.}

\subsubsection{Motocicletes fora de carretera}
La categoria fora de carretera (*off-road*) es divideix principalment en dues disciplines:

\begin{enumerate}
    \item \textbf{Motocròs:} Disciplina en circuit tancat amb salts i terreny irregular. Les màquines són molt lleugeres i no estan homologades per a la via pública.
    \item \textbf{Enduro:} Es practica en entorns naturals oberts. A diferència del motocròs, les motocicletes han d'estar homologades (llums, matrícula, miralls) per poder circular per camins legals.
\end{enumerate}

\end{document}